\section*{\texorpdfstring{\S\CurrentExerciseGroup}{Section \CurrentExerciseGroup}}
\phantomsection
\addcontentsline{toc}{section}{\texorpdfstring{Exercises for \S\CurrentExerciseGroup}{Exercises for Section \CurrentExerciseGroup}}

\begin{enumerate}
\def\labelenumi{\arabic{enumi})}
\item
  Let \(f\) and \(g\) be two measurable numerical functions, such that \(a = \mathrm{M}_{\infty}(f)\) and \(b = \mathrm{M}_{\infty}(g)\) are finite; show that, in order that \(\mathrm{M}_{\infty}(f + g) = \mathrm{M}_{\infty}(f) + \mathrm{M}_{\infty}(g)\) , it is necessary and sufficient that, for every pair of real numbers \(\alpha, \beta\) such that \(\alpha < a\) , \(\beta < b\) , the set of \(x \in X\) such that both \(\alpha \leqslant f(x)\) and \(\beta \leqslant g(x)\) not be locally negligible.
\item
  Let D be a closed convex set in a Banach space F, with nonempty interior, and let f be a measurable function with values in F, such that \(\mathbf{f}(\mathbf{X}) \subset \mathbf{D}\) . Let g be an integrable function \(\geqslant 0\) , not negligible, and such that fg is integrable. Suppose that the point
\end{enumerate}

\[
\mathbf {c} = \frac {\int \mathbf {f} g d \mu}{\int g d \mu}
\]

is a boundary point of \(\mathbf{D}\) ; show that if \(\mathbf{V}\) is the intersection of all the closed support hyperplanes of \(\mathbf{D}\) at the point \(\mathbf{c}\) , then \(\mathbf{f}(x) \in \mathbf{V} \cap \mathbf{D}\) almost everywhere in the set of points \(x\) such that \(g(x) > 0\) (reduce to the case that \(F\) is a separable space (that is, contains a countable dense set), using Th. 4 of §5; then note that \(V\) is the intersection of a countable family of support hyperplanes of \(D\) at the point \(c\) ).

Show that if F is finite-dimensional and if A is the facet of the point c with respect to D (TVS, II, §7, Exer. 3), the hypothesis implies that \(\mathbf{f}(x) \in \mathbf{A}\) almost everywhere in the set of points x such that \(g(x) > 0\) (argue by induction on the dimension of F).

\begin{enumerate}
\def\labelenumi{\arabic{enumi})}
\setcounter{enumi}{2}
\item
  Let \(\mu\) be a measure on a compact space \(X\) , such that \(X\) is equal to the support of \(\mu\) . Show that if every measurable numerical function that is bounded in measure is equal almost everywhere to a continuous function, then \(X\) is a Stone space (cf.~Ch. II, §1, Exer. 13) (consider the characteristic function of a compact set). Conversely, if \(X\) is a Stone space, in order that every measurable function that is bounded in measure be equal almost everywhere to a continuous function, it is necessary and sufficient that every nowhere dense set in \(X\) be negligible (cf.~§5, Exer. 18).
\item
  Let \(\varphi(t_1, t_2, \ldots, t_n)\) be a finite numerical function satisfying the conditions of Prop. 1 of Ch. I. Show that if \(f_1, f_2, \ldots, f_n\) are \(n\) finite numerical functions that are positive, integrable and non-negligible, then the function \(\varphi(f_1, f_2, \ldots, f_n)\) is integrable and
\end{enumerate}

\[
\int \varphi (f _ {1}, f _ {2}, \dots , f _ {n}) d \mu \leqslant \varphi (\int f _ {1} d \mu , \int f _ {2} d \mu , \dots , \int f _ {n} d \mu).
\]

Moreover, in order that

\[
\int \varphi (f _ {1}, \dots , f _ {n}) d \mu = \varphi (\int f _ {1} d \mu , \dots , \int f _ {n} d \mu),
\]

it is necessary and sufficient that, for almost every \(x \in X\) , the point of \(\mathbf{R}^n\) whose coordinates are \(\xi_i = f_i(x) / \varphi(f_1(x), \ldots, f_n(x))\) belong to the facet with respect to \(K\) of the point whose coordinates are

\[
\alpha_ {i} = (\int f _ {i} d \mu) / (\int \varphi (f _ {1}, \dots , f _ {n}) d \mu)
\]

(argue as in Exer. 2). In particular, \(p\) and \(q\) denoting two conjugate exponents such that \(1 < p < +\infty\) :

\(1^{\circ}\) for two positive numerical functions \(f \in \mathcal{L}^{p}\) , \(g \in \mathcal{L}^{q}\) to be such that \(\int fg d\mu = N_{p}(f)N_{q}(g)\) , it is necessary and sufficient that there exist two numbers \(\alpha, \beta\) , not both zero, such that \(\alpha(f(x))^{p} = \beta(g(x))^{q}\) almost everywhere;

\(2^{\circ}\) for two positive numerical functions \(f \in \mathcal{L}^{p}\) , \(g \in \mathcal{L}^{p}\) to be such that \(\mathrm{N}_{p}(f + g) = \mathrm{N}_{p}(f) + \mathrm{N}_{p}(g)\) , it is necessary and sufficient that there exist two numbers \(\alpha, \beta\) , not both zero, such that \(\alpha f(x) = \beta g(x)\) almost everywhere.

\begin{enumerate}
\def\labelenumi{\arabic{enumi})}
\setcounter{enumi}{4}
\tightlist
\item
  \begin{enumerate}
  \def\labelenumii{\alph{enumii})}
  \tightlist
  \item
    Let \(\mu\) be Lebesgue measure on the interval \(X = ]0, +\infty[\) . For every number \(p\) such that \(0 < p \leqslant +\infty\) , give examples of measurable functions \(f \geqslant 0\) on \(X\) , such that the set of numbers \(r\) ( \(0 \leqslant r \leqslant +\infty\) ), for which \(\mathrm{N}_r(f) = (\int f^r d\mu)^{1/r}\) is finite, is one of the intervals \([0, p[, ]0, p], ]p, +\infty]\) , \([p, +\infty]\) (take for \(f\) functions of the form \(x^\alpha(\log x)^\beta\) in the neighborhood of \(0\) or of \(+\infty\) ); from this, deduce that for every interval I contained in \([0, +\infty]\) , there exists a measurable function \(f \geqslant 0\) such that I is identical to the set of \(r > 0\) for which \(\mathrm{N}_r(f) < +\infty\) (consider the sum of two functions for which I has one of the above four forms).
  \end{enumerate}
\end{enumerate}

\begin{enumerate}
\def\labelenumi{\alph{enumi})}
\setcounter{enumi}{1}
\tightlist
\item
  For Lebesgue measure on the interval \(]0,1[\) , similarly give examples of functions f such that the set I of numbers r\textgreater0 for which \(\mathrm{N}_{r}(f)<+\infty\) is an arbitrary interval contained in \(]0,+\infty]\) with left end-point 0.
\end{enumerate}

\begin{enumerate}
\def\labelenumi{\arabic{enumi})}
\setcounter{enumi}{5}
\tightlist
\item
  For every measurable and non-negligible numerical function \(f \geqslant 0\) , show that \(\mathrm{N}_{r}(f)\) is an indefinitely differentiable function of r at every point interior to the interval where it is finite. From this, deduce that in the interior of the interval where \(\mathrm{N}_{r}(f)\)
\end{enumerate}

is finite, \(\log \mathrm{N}_r(f)\) is a strictly convex function of \(1 / r\) provided that \(f\) is not almost everywhere constant in the set of \(x\in \mathrm{X}\) where \(f(x)\neq 0\) .

¶ 7) Let \(f\) be a numerical function that is positive, measurable and non-negligible. a) Show that if, for a finite number \(r > 0\) , \(f^r\) is integrable, then \(\log f\) is quasi-integrable (§5, Exer. 15).

\begin{enumerate}
\def\labelenumi{\alph{enumi})}
\setcounter{enumi}{1}
\item
  Suppose that \(f^r\) is integrable for \(0 < r < r_0\) . Let \(A\) be the set of \(x \in X\) such that \(f(x) > 0\) ; show that if \(\mu^*(A) > 1\) , then \(N_r(f)\) tends to \(+\infty\) as \(r\) tends to \(0\) ; if \(\mu(A) < 1\) , \(N_r(f)\) tends to \(0\) with \(r\) (make use of Prop. 4 of Ch. I).
\item
  If \(\mu(A) = 1\) , show that \(\int f^r d\mu\) tends to 1 as \(r\) tends to 0, and has a right-derivative at this point, equal to \(\int \log f d\mu\) (make use of Exer. 12 of §5); from this, deduce that as \(r\) tends to 0, \(\mathrm{N}_r(f)\) tends to \(\mathrm{G}(f) = \exp(\int \log f d\mu)\) .
\item
  If \(\mu(\mathrm{X})=1\) and if \(f^{r}\) and \(\log f\) are integrable, show that \(\mathrm{G}(f)\leqslant\mathrm{N}_{r}(f)\) , the equality holding only if f is constant almost everywhere (make use of Exer. 4).
\item
  If \(\mu(X) = 1\) and if \(f\) and \(g\) are two measurable functions, positive and such that \(G(f)\) and \(G(g)\) are defined, show that \(G(f + g)\) is defined and that \(G(f) + G(g) \leqslant G(f + g)\) , the equality holding only if there exist two numbers \(\alpha, \beta\) , not both zero, such that \(\alpha f(x) = \beta g(x)\) almost everywhere, or if \(G(f + g) = 0\) (make use of \(d\) ), considering the functions \(f/(f + g)\) and \(g/(f + g)\) ).
\end{enumerate}

\begin{enumerate}
\def\labelenumi{\arabic{enumi})}
\setcounter{enumi}{7}
\tightlist
\item
  Let \(\mu\) be Lebesgue measure on the interval \(X = [0, +\infty[\) .
\end{enumerate}

\begin{enumerate}
\def\labelenumi{\alph{enumi})}
\item
  Let \(k > 1\) , \(h < k\) be two real numbers. Show that for every \(n \geqslant 1\) , and every number \(p\) such that \(1 \leqslant p \leqslant +\infty\) , the functions \(f_n(x) = n^h / (x + n)^k\) belong to \(\mathcal{L}^p\) ; show that \(\mathrm{N}_p(f_n)\) tends to 0 with \(1/n\) for \(p > 1 / (k - h)\) , but that the sequence of the \(\mathrm{N}_p(f_n)\) is not bounded for \(p < 1 / (k - h)\) .
\item
  Let \(k\) be a number \(< 1\) ; show that for every \(n > 1\) and every number \(p\) such that \(1 \leqslant p \leqslant +\infty\) , the functions \(g_n(x) = n^k e^{-nx}\) belong to \(\mathcal{L}^p\) ; show that \(\mathrm{N}_p(g_n)\) tends to 0 with \(1/n\) for \(p < 1/k\) , but that the sequence of the \(\mathrm{N}_p(g_n)\) is not bounded for \(p > 1/k\) .
\end{enumerate}

Deduce from \(a)\) and \(b)\) that if \(1 \leqslant p < q \leqslant +\infty\) , the topologies induced on \(\mathcal{L}^p \cap \mathcal{L}^q\) by those of \(\mathcal{L}^p\) and \(\mathcal{L}^q\) are not comparable.

\begin{enumerate}
\def\labelenumi{\alph{enumi})}
\setcounter{enumi}{2}
\tightlist
\item
  If \(\mu\) is Lebesgue measure on the interval \([0,1]\) , show similarly that if \(p < q\) , the topology of convergence in mean of order \(q\) is strictly finer than the topology of convergence in mean of order \(p\) (on \(\mathcal{L}^q\) ).
\end{enumerate}

\begin{enumerate}
\def\labelenumi{\arabic{enumi})}
\setcounter{enumi}{8}
\tightlist
\item
  Let \(X\) be an infinite discrete space, \(\mu\) a measure on \(X\) such that the support of \(\mu\) is equal to \(X\) .
\end{enumerate}

\begin{enumerate}
\def\labelenumi{\alph{enumi})}
\item
  Show that for \(1 \leqslant p \leqslant +\infty\) , the space \(\mathcal{L}^p(\mu)\) is a topological vector space isomorphic to the space \(\mathcal{L}^p(\mu_0)\) , where \(\mu_0\) is the measure on X defined by the mass \(+1\) at each point of X.
\item
  Show that if \(1 \leqslant p < q \leqslant +\infty\) , the topology of convergence in mean of order \(p\) is strictly finer than the topology of convergence in mean of order \(q\) (on \(\mathcal{L}^p\) ).
\end{enumerate}

¶ 10) a) In a Banach space F, let a and b be two vectors such that \(|a| = |b| = 1\) . Show that for every number t such that \(0 \leqslant t \leqslant 1\) , and for every p such that \(1 \leqslant p < +\infty\) ,

\[
| \mathbf {a} - t \mathbf {b} | ^ {p} \leqslant 2 ^ {p} | \mathbf {a} - t ^ {p} \mathbf {b} |\tag{1}
\]

\[
| \mathbf {a} - t ^ {p} \mathbf {b} | \leqslant 3 p | \mathbf {a} - t \mathbf {b} |\tag{2}
\]

(express \(\mathbf{a} - t^{p}\mathbf{b}\) as a linear combination of \(\mathbf{a} - tb\) and \(\mathbf{a} - \mathbf{b}\) and note that for \(0\leqslant p\leqslant 1\) , one has \(|\mathbf{a} - pb|\geqslant 1 - p\) and \(|\mathbf{a} - \mathbf{b}|\leqslant 2|\mathbf{a} - pb|)\) . From this, deduce that if \(\mathbf{y}\) and \(\mathbf{z}\) are any two vectors in \(\mathbb{F}\) , then

\[
\left| \mathbf {y} - \mathbf {z} \right| ^ {p} \leqslant 2 ^ {p} \left| \left| \mathbf {y} \right| ^ {p - 1} \cdot \mathbf {y} - \left| \mathbf {z} \right| ^ {p - 1} \cdot \mathbf {z} \right|\tag{3}
\]

\[
\left| | \mathbf {y} | ^ {p - 1} \cdot \mathbf {y} - | \mathbf {z} | ^ {p - 1} \cdot \mathbf {z} \right| \leqslant 3 p | \mathbf {y} - \mathbf {z} | (| \mathbf {y} | + | \mathbf {z} |) ^ {p - 1}.\tag{4}
\]

\begin{enumerate}
\def\labelenumi{\alph{enumi})}
\setcounter{enumi}{1}
\item
  Show that the mapping \(\mathbf{f} \mapsto |\mathbf{f}|^{(1/p)-1} \cdot \mathbf{f}\) is a uniformly continuous bijective mapping of \(\mathcal{L}_{\mathrm{F}}^{1}\) onto \(\mathcal{L}_{\mathrm{F}}^{p}\) (make use of the inequality (3)).
\item
  Show that the mapping \(\mathbf{f} \mapsto |\mathbf{f}|^{p-1} \cdot \mathbf{f}\) of \(\mathcal{L}_{\mathrm{F}}^{p}\) onto \(\mathcal{L}_{\mathrm{F}}^{1}\) is uniformly continuous on every bounded subset of the space \(\mathcal{L}_{\mathrm{F}}^{p}\) (make use of the inequality (4) and Hölder's inequality). Deduce from this that the topological spaces \(\mathcal{L}_{\mathrm{F}}^{1}\) and \(\mathcal{L}_{\mathrm{F}}^{p}\) are homeomorphic.
\end{enumerate}

\begin{enumerate}
\def\labelenumi{\arabic{enumi})}
\setcounter{enumi}{10}
\tightlist
\item
  Let \(g\) be a numerical function \(\geqslant 0\) belonging to \(\mathcal{L}^p\) ( \(1 \leqslant p < +\infty\) ). Denote by \(I_g\) the set of functions \(f \in \mathcal{L}_F^p\) such that \(|f| \leqslant g\) .
\end{enumerate}

\begin{enumerate}
\def\labelenumi{\alph{enumi})}
\item
  Show that on \(I_{g}\) , the topology of convergence in mean of order p is identical to the topology of convergence in measure.
\item
  If \(p \leqslant q < r\) (resp. \(q < r \leqslant p\) ), show that on the set \(I_g \cap \mathcal{L}_F^q \cap \mathcal{L}_F^r\) the topology of convergence in mean of order \(q\) is coarser (resp. finer) than the topology of convergence in mean of order \(r\) (for two functions \(f, f_0\) belonging to this set, write \(|f - f_0|^q \leqslant |f - f_0|^s(2g)^{q-s}\) and use Hölder's inequality, suitably choosing \(s\) and the pair of conjugate exponents). Show by means of examples that these topologies can be distinct (cf.~Exer. 8).
\item
  Let \(\mu\) be Lebesgue measure on \(X = ]0, +\infty[\); the function
\end{enumerate}

\[
g (x) = \left(x (\log^ {2} x + 1)\right) ^ {- 1 / p}
\]

belongs to \(\mathcal{L}^p\) , but to no \(\mathcal{L}^q\) for \(q \neq p\) . Show that if \(q \neq p\) , the topology of convergence in mean of order \(q\) on the set \(\mathrm{I}_g \cap \mathcal{L}^q\) is distinct from the topology of convergence in measure.

\begin{enumerate}
\def\labelenumi{\arabic{enumi})}
\setcounter{enumi}{11}
\tightlist
\item
  Let \(h\) be a numerical function \(\geqslant 0\) such that \(h\) and \(h^2\) are integrable; let \(I_h\) be the set of measurable numerical functions \(f\) such that \(|f| \leqslant h\) . Show that the mapping \((f, g) \mapsto fg\) of \(I_h \times I_h\) into \(\mathcal{L}^1\) is continuous for the topology of convergence in mean (on \(I_h\) and on \(\mathcal{L}^1\) ).
\end{enumerate}

¶ 13) For every number \(p\) such that \(0 < p < 1\) , denote by \(\mathcal{L}_F^p\) the set of measurable mappings \(\mathbf{f}\) of \(X\) into a Banach space \(F\) such that \(N sub p(\mathbf{f}) < +\infty\) .

\begin{enumerate}
\def\labelenumi{\alph{enumi})}
\item
  Show that \(\mathcal{L}_{\mathrm{F}}^{p}\) is a vector space and that, if \(\mathrm{B}_a\) denotes the set of \(\mathbf{f} \in \mathcal{L}_{\mathrm{F}}^{p}\) such that \(\mathrm{N}_p(\mathbf{f}) \leqslant a\) , the sets \(\mathrm{B}_a\) form, as \(a\) runs over the set of numbers \(>0\) , a fundamental system of neighborhoods of 0 for a metrizable topology compatible with the vector space structure of \(\mathcal{L}_{\mathrm{F}}^{p}\) .
\item
  Show that the mapping \(\mathbf{f} \mapsto |\mathbf{f}|^{p-1} \cdot \mathbf{f}\) is a uniformly continuous mapping of \(\mathcal{L}_{\mathrm{F}}^{p}\) onto \(\mathcal{L}_{\mathrm{F}}^{1}\) and that the inverse mapping is uniformly continuous on every bounded subset of \(\mathcal{L}_{\mathrm{F}}^{1}\) (cf.~Exer. 10). Deduce from this that the space \(\mathcal{L}_{\mathrm{F}}^{p}\) is complete, and that \(\mathcal{H}_{\mathrm{F}}(\mathrm{X})\) is dense in \(\mathcal{L}_{\mathrm{F}}^{p}\) .
\item
  If the measure \(\mu\) is bounded, then \(\mathcal{L}_{\mathrm{F}}^{1} \subset \mathcal{L}_{\mathrm{F}}^{p}\) and the topology of convergence in mean is finer than the topology induced on \(\mathcal{L}_{\mathrm{F}}^{1}\) by that of \(\mathcal{L}_{\mathrm{F}}^{p}\) .
\item
  Take \(\mu\) to be Lebesgue measure on \(X = [0,1]\) . Show that for every continuous function \(f \geqslant 0\) , there exists a decomposition \(f = \frac{1}{2}(f_1 + f_2)\) , where \(f_1\) and \(f_2\) are two functions \(\geqslant 0\) of \(\mathcal{L}^p\) such that \(\mathrm{N}_p(f_1) = \mathrm{N}_p(f_2) = 2^{1 - (1/p)}\mathrm{N}_p(f)\) . Deduce from this that, in \(\mathcal{L}^p\) , the closed convex envelope of every neighborhood \(\mathrm{B}_a\) is the entire space \(\mathcal{L}^p\) , consequently every continuous linear form on \(\mathcal{L}^p\) is identically zero.
\end{enumerate}

¶ 14) Let p and q be any two finite real numbers \textgreater{} 0 and let \(f(x_{1}, x_{2}, \ldots, x_{n})\) be a continuous numerical function defined on \(R^{n}\) .

\begin{enumerate}
\def\labelenumi{\alph{enumi})}
\tightlist
\item
  Let \(\mu\) be Lebesgue measure on \(X = [0,1]\) . In order that, for every system of \(n\) functions \(g_k \in \mathcal{L}^p\) , the function \(f(g_1, g_2, \ldots, g_n)\) belong to \(\mathcal{L}^q\) , it is necessary and sufficient that there exist a number \(a > 0\) such that
\end{enumerate}

\[
\left| f \left(x _ {1}, \dots , x _ {n}\right) \right| ^ {q} \leqslant a \left(1 + \left| x _ {1} \right| + \dots + \left| x _ {n} \right|\right) ^ {p}.
\]

(To see that the condition is necessary, argue by contradiction, supposing that for every integer \(m > 0\) , there exists a point \((x_{1m},\ldots ,x_{nm})\) of \(\mathbf{R}^n\) such that

\[
| f (x _ {1 m}, \dots , x _ {n m}) | ^ {q} \geqslant m (1 + | x _ {1 m} | + \dots + | x _ {n m} |) ^ {p}.
\]

Show that there would then exist in X a sequence \((\mathrm{A}_{m})\) of pairwise disjoint intervals such that, on setting \(g_{k}(t)=x_{km}\) for every \(t\in A_{m}\) and \(g_{k}(t)=0\) for every point t that does not belong to any of the \(A_{m}\) , each of the functions \(g_{k}\) would belong to \(L^{p}\) but \(f(g_{1},\ldots,g_{n})\) would not belong to \(L^{q}\) .

\begin{enumerate}
\def\labelenumi{\alph{enumi})}
\setcounter{enumi}{1}
\tightlist
\item
  Let \(\mu\) be Lebesgue measure on \(\mathbf{R}\) . In order that, for every system of \(n\) functions \(g_k \in \mathcal{L}^p\) , the function \(f(g_1, \ldots, g_n)\) belong to \(\mathcal{L}^q\) , it is necessary and sufficient that there exist a number \(b > 0\) such that
\end{enumerate}

\[
\left| f \left(x _ {1}, \dots , x _ {n}\right) \right| ^ {q} \leqslant b \left(\left| x _ {1} \right| + \left| x _ {2} \right| + \dots + \left| x _ {n} \right|\right) ^ {p}
\]

(same method).

¶ 15) Let X be a compact space, \(\mu\) a positive measure on X. For two functions \(f, g\) of \(\mathcal{L}_{\mathbf{C}}^{2}\) , one sets \((f|g) = (\widetilde{f}|\widetilde{g}) = \int f\overline{g} d\mu\) . A sequence of functions \(f_{n} \in \mathcal{L}_{\mathbf{C}}^{2}\) is said to be orthonormal if the sequence \(\widetilde{f}_{n}\) is orthonormal in the Hilbert space \(\mathrm{L}_{\mathbf{C}}^{2}\) , that is (TVS, V, §2) if \((f_{m}|f_{n}) = \delta_{mn}\) (the Kronecker symbol) for every pair of indices. For every function \(g \in \mathcal{L}_{\mathbf{C}}^{2}\) , the complex numbers \(c_{n} = (g|f_{n})\) are called the components of \(g\) with respect to the orthonormal sequence \((f_{n})\) ; one has \(\sum_{n=0}^{\infty} |c_{n}|^{2} \leqslant \int |g|^{2} d\mu\) .

For every pair of points \(x, y\) of \(X\) and every integer \(n \geqslant 0\) , one sets \(K_n(x, y) = \sum_{k=0}^{n} f_k(x) \overline{f_k(y)}\) (the \(n\) -th kernel of the orthonormal sequence \((f_n)\) ); for every function \(g \in \mathcal{L}_{\mathbf{C}}^2\) ,

\[
s _ {n} (g) = \sum_ {k = 0} ^ {n} (g | f _ {k}) f _ {k} (x) = \int \mathrm{K} _ {n} (x, y) g (y) d \mu (y).
\]

Set \(\mathrm{H}_n(x) = \int |\mathrm{K}_n(x,y)|d\mu (y)\) (called the \(n\) -th Lebesgue function of the orthonormal sequence \((f_n)\) ).

\begin{enumerate}
\def\labelenumi{\alph{enumi})}
\item
  Let \((\alpha_{n})\) be a decreasing sequence of numbers \(>0\) such that the series with general term \(\alpha_{n}\) is convergent. Show that, for almost every \(x \in X\) , \(\sum_{k=0}^{n} |f_{k}(x)|^{2} = o(1/\alpha_{n})\) (using Prop. 6 of §3, show that the series with general term \(\alpha_{n}|f_{n}(x)|^{2}\) is convergent almost everywhere, and make use of Exer. 10 of GT, IV, §7). Deduce from this that \(\mathrm{H}_{n}(x) = o(1/\sqrt{\alpha_{n}})\) for almost every \(x \in X\) .
\item
  Let \(x_0\) be a point of \(X\) . In order that, for every complex-valued function \(g\) , defined and continuous on \(X\) , the partial sums \(s_n(g)\) be bounded at the point \(x_0\) (by a number depending on \(g\) and \(x_0\) ), it is necessary and sufficient that the set of numbers \(H_n(x_0)\) be bounded (make use of Prop. 3 and the fact that in the dual of a Banach space, every weakly bounded set is strongly bounded).
\item
  In order that, for every complex-valued function \(g\) , defined and continuous on \(X\) , the series with general term \((g|f_n)f_n(x)\) be uniformly convergent in \(X\) and have sum \(g(x)\) , it is necessary and sufficient that: \(1^\circ\) every continuous complex-valued function on \(X\) be uniformly approximable by linear combinations of the \(f_k\) ; \(2^\circ\) there exist a constant \(a\) such that \(|\mathrm{H}_n(x)| \leqslant a\) for all \(n\) and \(x \in X\) . (Note that for every \(n\) , \(f_n(x) = \sum_{m=0}^{\infty}(f_n|f_m)f_m(x)\) identically; on the other hand, to prove the necessity of the condition \(2^{\circ}\) , note that for every increasing sequence \((n_{k})\) of integers and every sequence \((x_{k})\) of points of X, the sequence of numbers \(\int \mathrm{K}_{n_k}(x_k,y)g(y)d\mu (y)\) is bounded (by a number depending on \(g\) ) and argue as in \(b\) .)
\end{enumerate}

\begin{enumerate}
\def\labelenumi{\arabic{enumi})}
\setcounter{enumi}{15}
\tightlist
\item
  Let \(\mu\) be Lebesgue measure on \(X = [0, 2\pi]\) . The sequence of functions \(f_n\) such that
\end{enumerate}

\[
f _ {0} (x) = \frac {1}{\sqrt {2 \pi}}, \quad f _ {2 n - 1} (x) = \frac {1}{\sqrt {\pi}} \cos n x, \quad f _ {2 n} (x) = \frac {1}{\sqrt {\pi}} \sin n x \quad (n \geqslant 1)
\]

is an orthonormal and total sequence in the space \(\mathcal{L}_{\mathbf{C}}^{2}\) (cf.~GT, X, §4, Prop. 8). Show that the corresponding Lebesgue function \(\mathrm{H}_n(x)\) is independent of \(x\) and that \(\mathrm{H}_n \sim 4 / \pi \log n\) .

¶ 17) Let \(\mu\) be Lebesgue measure on \(X = [0,1]\) . One defines the sequence \((f_n)\) of step functions on \(X\) by the following conditions: \(f_0\) is the constant 1; for every integer \(n > 0\) , let \(m\) be the largest integer such that \(2^m \leqslant n\) , and write \(n = 2^m + k\) ; \(f_n\) is the function equal to \(2^{m/2}\) on the interval \(\left[\frac{2k}{2^{m+1}}, \frac{2k+1}{2^{m+1}}\right]\) , to \(-2^{m/2}\) on the interval \(\left[\frac{2k+1}{2^{m+1}}, \frac{2k+2}{2^{m+1}}\right]\) , and to 0 at the other points of \(X\) .

\begin{enumerate}
\def\labelenumi{\alph{enumi})}
\item
  Show that the sequence \((f_n)\) is orthonormal (the Haar orthonormal system).
\item
  Let \(V_{n}\) be the linear subspace of \(L_{C}^{2}\) (over C) generated by the \(f_{k}\) with indices \(k \leqslant n\) . Show that there exists a partition of {[}0,1{[} into \(n+1\) semi-open intervals such that, in each of these intervals, every function belonging to \(V_{n}\) is constant. Deduce from this that, conversely, for every function g that is constant on each of these \(n+1\) intervals, there exists a function in \(V_{n}\) that is equal to g on {[}0,1{[} (note that \(V_{n}\) has dimension \(n+1\) ).
\item
  Let \(g\) be any function in \(\mathcal{L}_{\mathbf{C}}^2\) ; deduce from \(b\) ) that if \(h\) is the unique function in \(V_n\) for which \(N_2(g - h)\) is minimal, then, in every interval \([\alpha, \beta[\) where \(h\) is constant, one has \(h(x) = \frac{1}{\beta - \alpha} \int_{\alpha}^{\beta} g(t) dt\) .
\item
  Show that for every complex-valued function g, defined and continuous on X, the series with general term \((g|f_{n})f_{n}(x)\) is uniformly convergent on {[}0,1{[} and has sum \(g(x)\) (make use of c)). Deduce from this that the sequence \((f_{n})\) is total.
\end{enumerate}

¶ 18) Let X be a locally compact space, μ a positive measure on X, and f a μ-measurable numerical function.

\begin{enumerate}
\def\labelenumi{\alph{enumi})}
\tightlist
\item
  Let \((a_{n})_{n\in \mathbf{Z}}\) be a sequence of real numbers, \(\lambda\) a finite real number; for every \(n\in \mathbf{Z}\) , set
\end{enumerate}

\[
u _ {n} (x) = a _ {n + [ \lambda \log | f (x) | ]} f (x) \quad \text { if } f (x) \neq 0 \text { and } f (x) \neq \pm \infty
\]

\[
u _ {n} (x) = 0 \quad \text {   for   the   other   values   of   } x \in X
\]

(where \([t]\) denotes the integral part of the finite real number \(t\) ). Show that the \(u_{n}\) are \(\mu\) -measurable and that, for every finite real number \(c\) such that \(c\lambda < 1\) ,

\[
\sum_ {n = - \infty} ^ {+ \infty} \int^ {*} | u _ {n} (x) | e ^ {c x} d \mu (x) \leqslant e ^ {| c |} \int^ {*} | f (x) | ^ {1 - c \lambda} d \mu (x) \cdot \left(\sum_ {n = - \infty} ^ {+ \infty} e ^ {c n} | a _ {n} |\right).
\]

\begin{enumerate}
\def\labelenumi{\alph{enumi})}
\setcounter{enumi}{1}
\tightlist
\item
  Let \(p', p''\) be two finite numbers \(> 0\) , \(t\) a real number such that \(0 < t < 1\) , and let \(p\) be the real number defined by \(1/p = (1 - t)/p' + t/p''\) . For every number \(\alpha > 0\) , set
\end{enumerate}

\[
1 / \mathrm{K} _ {\alpha} = \inf \left(\sum_ {n = - \infty} ^ {+ \infty} | e ^ {\alpha t n} a _ {n} | ^ {p ^ {\prime}}\right) ^ {(1 - t) / p ^ {\prime}} \left(\sum_ {n = - \infty} ^ {+ \infty} | e ^ {- \alpha (1 - t) n} a _ {n} | ^ {p ^ {\prime \prime}}\right) ^ {t / p ^ {\prime \prime}},
\]

the infimum being taken over all the absolutely convergent series \((a_{n})_{n\in \mathbf{Z}}\) of finite real numbers such that \(\sum_{n = -\infty}^{+\infty}a_n = 1\) . Show that

(*)

\[
\begin{array}{l} \mathrm{N} _ {p} (f) \leqslant \mathrm{K} _ {\alpha} \cdot \inf \mathrm{F} (u) \\ \mathrm{N} _ {p} (f) \geqslant \mathrm{K} _ {\alpha} e ^ {- \alpha} \cdot \inf \mathrm{F} (u), \end{array}\tag{**}
\]

where \(u = (u_n)_{n \in \mathbf{Z}}\) runs over the set of sequence of functions belonging to \(\mathcal{L}^{p'} \cap \mathcal{L}^{p''}\) such that the series \((u_n(x))_{n \in \mathbf{Z}}\) is almost everywhere absolutely convergent with sum \(f(x)\) ; for every sequence \(u\) having these properties, one sets

\[
\mathrm{F} (u) = \left(\sum_ {n = - \infty} ^ {+ \infty} \left(\mathrm{N} _ {p ^ {\prime}} (e ^ {\alpha t n} u _ {n})\right) ^ {p ^ {\prime}}\right) ^ {(1 - t) / p ^ {\prime}} \left(\sum_ {n = - \infty} ^ {+ \infty} \left(\mathrm{N} _ {p ^ {\prime \prime}} (e ^ {- \alpha (1 - t) n} u _ {n})\right) ^ {p ^ {\prime \prime}}\right) ^ {t / p ^ {\prime \prime}}
\]

and in the formulas (*) and (**) the infimum is taken over the set of sequences \(u = (u_n)\) having the preceding properties. (To prove (*), use Hölder's inequality; to prove (**), use a) twice with suitable choices of \(c\) and \(\lambda\) .) In particular, \(K_{\alpha}\) is finite for all \(\alpha > 0\) .

\begin{enumerate}
\def\labelenumi{\alph{enumi})}
\setcounter{enumi}{2}
\tightlist
\item
  Let \(p', p'', q', q''\) be finite numbers \(\geqslant 1\) , \(t\) a number such that \(0 < t < 1\) , and
\end{enumerate}

let

\[
{\frac {1}{p}} = {\frac {1 - t}{p ^ {\prime}}} + {\frac {t}{p ^ {\prime \prime}}}, \quad {\frac {1}{q}} = {\frac {1 - t}{q ^ {\prime}}} + {\frac {t}{q ^ {\prime \prime}}}.
\]

Let Y be a second locally compact space, \(\nu\) a positive measure on Y, and let w be a linear mapping of \(\mathcal{K}(\mathrm{X};\mathbf{R})\) into the vector space (not topological) \(\mathcal{S}(\mathrm{Y},\nu;\mathbf{R})\) of \(\nu\) -measurable finite numerical functions on Y. Suppose that: \(1^{\circ}w\) maps \(\mathcal{K}(\mathrm{X};\mathbf{R})\) into \(\mathcal{L}^{q'}(\mathrm{Y};\mathbf{R})\cap\mathcal{L}^{q''}(\mathrm{Y};\mathbf{R})\) ; \(2^{\circ}\) for every function \(f\in\mathcal{K}(\mathrm{X};\mathbf{R})\) ,

\[
\mathrm{N} _ {q ^ {\prime}} (w (f)) \leqslant \mathrm{M} ^ {\prime} \mathrm{N} _ {p ^ {\prime}} (f) \quad \text { and } \quad \mathrm{N} _ {q ^ {\prime \prime}} (w (f)) \leqslant \mathrm{M} ^ {\prime \prime} \mathrm{N} _ {p ^ {\prime \prime}} (f).
\]

From this, conclude that \(w\) also maps \(\mathcal{K}(\mathrm{X};\mathbf{R})\) into \(\mathcal{L}^q (\mathrm{Y};\mathbf{R})\) and that

\[
\mathrm{N} _ {q} (w (f)) \leqslant \mathrm{M} \cdot \mathrm{N} _ {p} (f)
\]

with

\[
\mathrm{M} \leqslant \mathrm{M} ^ {\prime 1 - t} \mathrm{M} ^ {\prime \prime t}
\]

(M. Riesz's inequality). (Write \(f\) in the form \(\sum_{n} u_n\) as in \(b\) ), and consider the series with general term \(w(u_n)\) ; use the inequalities (*) and (** of \(b\) ).)

\begin{enumerate}
\def\labelenumi{\arabic{enumi})}
\setcounter{enumi}{18}
\tightlist
\item
  Let \(f\) be a numerical function \(\geqslant 0\) defined in the space \(\mathbf{R}_+^* = ]0, +\infty[\) and \(p\) -th power integrable for Lebesgue measure \((1 < p < +\infty)\) . Set \(F(x) = \int_0^x f(t) dt\) for all \(x > 0\) .
\end{enumerate}

\begin{enumerate}
\def\labelenumi{\alph{enumi})}
\item
  Show that as \(x\) tends to 0 or to \(+\infty\) , \(F(x) = o(x^{(p - 1) / p})\) (use Hölder's inequality).
\item
  Show that the function \(\mathrm{F}(x)/x\) is p-th power integrable in \(R_{+}^{*}\) and that
\end{enumerate}

\[
\int_ {0} ^ {+ \infty} \left(\frac {\mathrm{F} (x)}{x}\right) ^ {p} d x \leqslant \left(\frac {p}{p - 1}\right) ^ {p} \int_ {0} ^ {+ \infty} (f (x)) ^ {p} d x
\]

(Hardy's inequality). (Consider first the case that \(f \in \mathcal{K}(\mathbf{R}_+^*)\) ; for every compact interval \([a, b] \subset \mathbf{R}_+^*\) , obtain an upper bound for the integral

\[
\int_ {a} ^ {b} \left(\frac {\mathrm{F} (t)}{t}\right) ^ {p} d t
\]

by integrating by parts and making use of Hölder's inequality.)

¶ 20) a) Let Y be a metric space; for every \(y \in Y\) and every \(r > 0\) , denote by B(y;r) the open ball in Y with center y and radius r. Let \(\mathfrak{S}\) be a set of open balls in Y whose diameters form a bounded set in R, and which is such that, for every sequence (B(yn;r\_n)) of pairwise disjoint balls belonging to \(\mathfrak{S}\) , one has \(\lim_{n \to \infty} r_n = 0\) . Show that if M is the union of the balls B \(\in \mathfrak{S}\) , there exists a sequence of pairwise disjoint balls B(yn;r\_n) \(\in \mathfrak{S}\) such that the balls B(yn;4r\_n) form a covering of M. (If k \textgreater{} 0 is an upper bound for the set of radii of the balls B \(\in \mathfrak{S}\) , define by induction on h a sequence of families ( \(\mathfrak{F}_h\) ) of balls B(yhj;r\_hj) \(\in \mathfrak{S}\) such that \(\mathfrak{F}_h\) is maximal among the (finite) families of balls belonging to \(\mathfrak{S}\) , pairwise disjoint and disjoint from the balls belonging to the families \(\mathfrak{F}_i\) for \(i < h\) , and with radii between \((2/3)^{h+1}k\) and \((2/3)^h k\) .)

\begin{enumerate}
\def\labelenumi{\alph{enumi})}
\setcounter{enumi}{1}
\tightlist
\item
  Let X be a metrizable locally compact space, d a metric on X compatible with its topology; again denote by \(\mathrm{B}(x;r)\) the open ball with center x and radius r \textgreater{} 0 for this metric, and by \(\delta(\mathrm{A})\) the diameter of a subset A of X for the metric d.~Let \(\mu\) be a positive measure on X satisfying the following conditions: \(1^{\circ}\) every open ball is \(\mu\) -integrable; \(2^{\circ}\) \(\mu(\mathrm{B}(x;4r)) \leqslant \mathrm{K} \cdot \mu(\mathrm{B}(x;r))\) , where K is a constant \textgreater{} 1 independent of x and r; \(3^{\circ}\) if a sequence \((\mathrm{B}_{n})\) of open balls is such that \(\lim_{n \to \infty} \mu(\mathrm{B}_{n}) = 0\) , then \(\lim_{n \to \infty} \delta(\mathrm{B}_{n}) = 0\) ; \(4^{\circ}\) if a sequence \((\mathrm{B}_{n})\) of open balls is such that \(\lim_{n \to \infty} \delta(\mathrm{B}_{n}) = +\infty\) , then \(\lim_{n \to \infty} \mu(\mathrm{B}_{n}) = +\infty\) .
\end{enumerate}

Let \(f \in \mathcal{L}^p(X; \mu)\) be a function with values \(\geqslant 0\) ( \(1 < p < +\infty\) ); for every \(x \in X\) , set

\[
\overline {{{{f}}}} (x) = \sup \frac {1}{\mu (\mathrm{B})} \int_ {\mathrm{B}} f (y) d \mu (y),
\]

where B runs over the set of open balls with center \(x\) . Show that \(\overline{f}\) is finite and lower semi-continuous on \(X\) .

\begin{enumerate}
\def\labelenumi{\alph{enumi})}
\setcounter{enumi}{2}
\tightlist
\item
  Let \(f^*\) be the decreasing rearrangement of \(f\) (§5, Exer. 29); for \(t > 0\) , set
\end{enumerate}

\[
\beta_ {f} (t) = \frac {1}{t} \int_ {0} ^ {t} f ^ {*} (s) d s,
\]

which is a decreasing continuous function such that \(f^* \leqslant \beta_f\) ; one has \(\beta_f(t) = o(t^{-1/p})\) as \(t\) tends to 0 or to \(+\infty\) (Exer. 19 a)). Finally, denote by \(\gamma_f\) the inverse function of \(\beta_f\) , defined on the interval \(]0, \beta_f(0+)[\) , and extended by 0 to the exterior of this interval if \(\beta_f(0+)\) is finite.

For every \(t > 0\) , denote by \(\mathrm{M}_t\) the set of \(x \in X\) such that \(\overline{f}(x) > t\) . Show that \(\mu(\mathrm{M}_t) \leqslant \mathrm{K} \cdot \gamma_f(t)\) . (For every \(x \in \mathrm{M}_t\) , let \(\mathrm{B}_x\) be an open ball with center \(x\) such that \(\int_{\mathrm{B}_x} f(y) d\mu(y) \geqslant t \cdot \mu(\mathrm{B}_x)\) ; apply a) to the family of balls \(\mathrm{B}_x\) .)

\(d\) ) Deduce from \(c\) ) that

\[
\int \left(\overline {{f}} (x)\right) ^ {p} d \mu (x) \leqslant \mathrm{K} \left(\frac {p}{p - 1}\right) ^ {p} \int \left(f (x)\right) ^ {p} d \mu (x).
\]

(Note that the first member may also be written \(\int_0^{+\infty}pt^{p - 1}\big(\mu (\mathrm{M}_t)\big)dt\) , and make use of Hardy's inequality (Exer. 19).)
% semantic-fragments: legacy-input-seam
\relax{}
\setcounter{exercisegroup}{7}
\edef\CurrentExerciseGroup{\arabic{exercisegroup}}
